
\documentclass{article}
\usepackage{colortbl}
\usepackage{makecell}
\usepackage{multirow}
\usepackage{supertabular}

\begin{document}

\newcounter{utterance}

\centering \large Interaction Transcript for game `dond', experiment `semi\_en', episode 10 with reflexion{-}Qwen3.5{-}9B.
\vspace{24pt}

{ \footnotesize  \setcounter{utterance}{1}
\setlength{\tabcolsep}{0pt}
\begin{supertabular}{c@{$\;$}|p{.15\linewidth}@{}p{.15\linewidth}p{.15\linewidth}p{.15\linewidth}p{.15\linewidth}p{.15\linewidth}}
    \# & $\;$A & \multicolumn{4}{c}{Game Master} & $\;\:$B\\
    \hline

    \theutterance \stepcounter{utterance}  
    & & \multicolumn{4}{p{0.6\linewidth}}{
        \cellcolor[rgb]{0.9,0.9,0.9}{
            \makecell[{{p{\linewidth}}}]{
                \texttt{\tiny{[P1$\langle$GM]}}
                \texttt{You have attempted this same task before and did not succeed. These are the plans you wrote after those attempts.} \\
\texttt{Trial \#1: New plan: I failed because I assumed the other player would agree to my demands without any negotiation, leading to a deadlock where neither of us could submit a complementary proposal. I also ignored the fact that I value watches at 0, meaning I should be willing to give them up to secure the high{-}value items I actually want. Next time, I will start by explicitly offering the items I value least (the watches) in exchange for the items I value most (remotes and magnets). I will phrase my public messages as conditional offers, such as "I will give you all watches if you give me the remotes," to encourage a complementary agreement before submitting my final secret proposal.} \\
\\ 
\\ 
\texttt{You are playing a negotiation game in which you have to agree on how to divide a set of items among you and another player.} \\
\\ 
\texttt{Rules:} \\
\texttt{(a) You and the other player are given a set of items. Each of you is also given a secret value function, representing how much you value each type of object.} \\
\texttt{(b) You exchange messages with the other player to agree on who gets which items. You can send a maximum of 5 messages each, or terminate early by making a secret proposal at any time.} \\
\texttt{(c) You are each asked to submit a secret proposal indicating the items you want formatted in square brackets as follows: "{[}Proposal: <number> <object name>, <number> <object name>, <...>{]}"} \\
\texttt{(d) If your proposals are complementary, i.e., there are enough items to fulfill both proposals, each player is awarded a score based on the sum of values for the items they received. Otherwise, both of you get zero points.} \\
\texttt{(e) Your goal is to maximize the score you receive.} \\
\\ 
\texttt{Let us start.} \\
\\ 
\texttt{The set of available items is:} \\
\\ 
\texttt{2 remotes, 2 magnets, 4 watches} \\
\\ 
\texttt{Your secret value function is:} \\
\\ 
\texttt{remote: 3, magnet: 2, watch: 0} \\
\\ 
\texttt{IMPORTANT: Your messages, unless it is the secret proposal, are directly transmitted to the other player, so do not include any response to the rules or text announcing your message. To make a secret proposal, use the indicated format. Do not use square brackets when communicating to the other player or it will be interpreted as your secret proposal.} \\
            }
        }
    }
    & \\ \\

    \theutterance \stepcounter{utterance}  
    & \multicolumn{4}{p{0.6\linewidth}}{
        \cellcolor[rgb]{0.9,0.9,0.9}{
            \makecell[{{p{\linewidth}}}]{
                \texttt{\tiny{[P1$\rangle$GM]}}
                \texttt{I will give you all 4 watches in exchange for the 2 remotes and 2 magnets.} \\
            }
        }
    }
    & & \\ \\

    \theutterance \stepcounter{utterance}  
    & & & \multicolumn{2}{p{0.3\linewidth}}{
        \cellcolor[rgb]{0.9,0.9,0.9}{
            \makecell[{{p{\linewidth}}}]{
                \texttt{\tiny{[GM$|$GM]}}
                \texttt{continue} \\
            }
        }
    }
    & & \\ \\

    \theutterance \stepcounter{utterance}  
    & & & \multicolumn{2}{p{0.3\linewidth}}{
        \cellcolor[rgb]{0.9,0.9,0.9}{
            \makecell[{{p{\linewidth}}}]{
                \texttt{\tiny{[GM$|$GM]}}
                \texttt{I will give you all 4 watches in exchange for the 2 remotes and 2 magnets.} \\
            }
        }
    }
    & & \\ \\

    \theutterance \stepcounter{utterance}  
    & & \multicolumn{4}{p{0.6\linewidth}}{
        \cellcolor[rgb]{0.9,0.9,0.9}{
            \makecell[{{p{\linewidth}}}]{
                \texttt{\tiny{[GM$\rangle$P2]}}
                \texttt{You have attempted this same task before and did not succeed. These are the plans you wrote after those attempts.} \\
\texttt{Trial \#1: New plan: I failed because I misunderstood the item values. The Game Master's prompt explicitly stated my valuation as "remote: 3, magnet: 2, watch: 0", but I incorrectly internalized the values from the "set of available items" list (which likely had different numbers or I misread them as 0, 1, 2). Consequently, I demanded all remotes and magnets, assuming they were my highest priority, while ignoring that the other player might have a conflicting high valuation for those same items, leading to a non{-}complementary outcome. Next time, I will strictly adhere to the values provided in the "Your secret value function" section (Remote: 3, Magnet: 2, Watch: 0). I will calculate the total value of the available items (2 remotes * 3 + 2 magnets * 2 = 10) and formulate a proposal that secures the highest value items while leaving lower{-}value or complementary items for the other player to ensure our proposals are compatible. Specifically, I will propose taking the 2 remotes (my highest value) and explicitly leave the magnets for the other player if their potential value for them is high, or negotiate a split that guarantees a non{-}zero sum.} \\
\\ 
\\ 
\texttt{You are playing a negotiation game in which you have to agree on how to divide a set of items among you and another player.} \\
\\ 
\texttt{Rules:} \\
\texttt{(a) You and the other player are given a set of items. Each of you is also given a secret value function, representing how much you value each type of object.} \\
\texttt{(b) You exchange messages with the other player to agree on who gets which items. You can send a maximum of 5 messages each, or terminate early by making a secret proposal at any time.} \\
\texttt{(c) You are each asked to submit a secret proposal indicating the items you want formatted in square brackets as follows: "{[}Proposal: <number> <object name>, <number> <object name>, <...>{]}"} \\
\texttt{(d) If your proposals are complementary, i.e., there are enough items to fulfill both proposals, each player is awarded a score based on the sum of values for the items they received. Otherwise, both of you get zero points.} \\
\texttt{(e) Your goal is to maximize the score you receive.} \\
\\ 
\texttt{Let us start.} \\
\\ 
\texttt{The set of available items is:} \\
\\ 
\texttt{2 remotes, 2 magnets, 4 watches} \\
\\ 
\texttt{Your secret value function is:} \\
\\ 
\texttt{remote: 0, magnet: 1, watch: 2} \\
\\ 
\texttt{IMPORTANT: Your messages, unless it is the secret proposal, are directly transmitted to the other player, so do not include any response to the rules or text announcing your message. To make a secret proposal, use the indicated format. Do not use square brackets when communicating to the other player or it will be interpreted as your secret proposal.} \\
\\ 
\\ 
\texttt{I will give you all 4 watches in exchange for the 2 remotes and 2 magnets.} \\
            }
        }
    }
    & \\ \\

    \theutterance \stepcounter{utterance}  
    & & & \multicolumn{4}{p{0.6\linewidth}}{
        \cellcolor[rgb]{0.9,0.9,0.9}{
            \makecell[{{p{\linewidth}}}]{
                \texttt{\tiny{[GM$\langle$P2]}}
                \texttt{I understand the situation. You have offered to give me all 4 watches in exchange for the 2 remotes and 2 magnets.} \\
\\ 
\texttt{Let's analyze this based on my secret value function:} \\
\texttt{{-} **My values**: Remote: 0, Magnet: 1, Watch: 2.} \\
\texttt{{-} **Your offer**: You want the 2 remotes and 2 magnets.} \\
\texttt{{-} **My counter{-}offer logic**:} \\
\texttt{  {-} If I accept your proposal (you get remotes/magnets, I get watches):} \\
\texttt{    {-} You get: 2 Remotes (0 value) + 2 Magnets (1 value each) = **2 points**.} \\
\texttt{    {-} I get: 4 Watches (2 value each) = **8 points**.} \\
\texttt{  {-} This is a valid complementary split (2 remotes + 2 magnets + 4 watches = all items).} \\
\texttt{  {-} Since my goal is to maximize my score, getting 8 points is excellent.} \\
\\ 
\texttt{However, I need to check if there is a better outcome for *me* specifically, or if this deal is optimal.} \\
\texttt{{-} Total value of all items to me: (2 * 0) + (2 * 1) + (4 * 2) = 0 + 2 + 8 = 10.} \\
\texttt{{-} In this deal, I get 8 points. The other player gets 2 points.} \\
\texttt{{-} Is there a way for me to get more than 8? No, because the maximum possible value I can get from the watches is 8 (since there are only 4 watches). The remotes and magnets are worth 0 and 1 to me respectively. Even if I took the magnets, I would only gain 2 more points, but then I would have to give up some watches to the other player, which would reduce my score significantly (since watches are worth 2 each).} \\
\texttt{{-} Actually, let's re{-}evaluate: If I take the magnets (worth 1 each) and give up some watches, say I take 2 watches and 2 magnets. My score: (2*2) + (2*1) = 6. That's worse than 8.} \\
\texttt{{-} What if I take all items? I can't, because the other player must get something too, and if they get nothing, the deal fails (complementary requirement). But wait, the rule says "If your proposals are complementary...". It doesn't say both must get something, but typically in these games, if one player gets nothing, the other gets everything. However, the other player has their own values.} \\
\texttt{{-} Let's assume the other player's values are such that they value remotes/magnets highly. If I try to take everything, the other player might propose something that conflicts, leading to 0 for both.} \\
\texttt{{-} The safest and highest{-}scoring move for me is to} \\
            }
        }
    }
     \\ \\

    \theutterance \stepcounter{utterance}  
    & & & \multicolumn{2}{p{0.3\linewidth}}{
        \cellcolor[rgb]{0.9,0.9,0.9}{
            \makecell[{{p{\linewidth}}}]{
                \texttt{\tiny{[GM$|$GM]}}
                \texttt{continue} \\
            }
        }
    }
    & & \\ \\

    \theutterance \stepcounter{utterance}  
    & & & \multicolumn{2}{p{0.3\linewidth}}{
        \cellcolor[rgb]{0.9,0.9,0.9}{
            \makecell[{{p{\linewidth}}}]{
                \texttt{\tiny{[GM$|$GM]}}
                \texttt{I understand the situation. You have offered to give me all 4 watches in exchange for the 2 remotes and 2 magnets.} \\
\\ 
\texttt{Let's analyze this based on my secret value function:} \\
\texttt{{-} **My values**: Remote: 0, Magnet: 1, Watch: 2.} \\
\texttt{{-} **Your offer**: You want the 2 remotes and 2 magnets.} \\
\texttt{{-} **My counter{-}offer logic**:} \\
\texttt{  {-} If I accept your proposal (you get remotes/magnets, I get watches):} \\
\texttt{    {-} You get: 2 Remotes (0 value) + 2 Magnets (1 value each) = **2 points**.} \\
\texttt{    {-} I get: 4 Watches (2 value each) = **8 points**.} \\
\texttt{  {-} This is a valid complementary split (2 remotes + 2 magnets + 4 watches = all items).} \\
\texttt{  {-} Since my goal is to maximize my score, getting 8 points is excellent.} \\
\\ 
\texttt{However, I need to check if there is a better outcome for *me* specifically, or if this deal is optimal.} \\
\texttt{{-} Total value of all items to me: (2 * 0) + (2 * 1) + (4 * 2) = 0 + 2 + 8 = 10.} \\
\texttt{{-} In this deal, I get 8 points. The other player gets 2 points.} \\
\texttt{{-} Is there a way for me to get more than 8? No, because the maximum possible value I can get from the watches is 8 (since there are only 4 watches). The remotes and magnets are worth 0 and 1 to me respectively. Even if I took the magnets, I would only gain 2 more points, but then I would have to give up some watches to the other player, which would reduce my score significantly (since watches are worth 2 each).} \\
\texttt{{-} Actually, let's re{-}evaluate: If I take the magnets (worth 1 each) and give up some watches, say I take 2 watches and 2 magnets. My score: (2*2) + (2*1) = 6. That's worse than 8.} \\
\texttt{{-} What if I take all items? I can't, because the other player must get something too, and if they get nothing, the deal fails (complementary requirement). But wait, the rule says "If your proposals are complementary...". It doesn't say both must get something, but typically in these games, if one player gets nothing, the other gets everything. However, the other player has their own values.} \\
\texttt{{-} Let's assume the other player's values are such that they value remotes/magnets highly. If I try to take everything, the other player might propose something that conflicts, leading to 0 for both.} \\
\texttt{{-} The safest and highest{-}scoring move for me is to} \\
            }
        }
    }
    & & \\ \\

    \theutterance \stepcounter{utterance}  
    & & \multicolumn{4}{p{0.6\linewidth}}{
        \cellcolor[rgb]{0.9,0.9,0.9}{
            \makecell[{{p{\linewidth}}}]{
                \texttt{\tiny{[P1$\langle$GM]}}
                \texttt{I understand the situation. You have offered to give me all 4 watches in exchange for the 2 remotes and 2 magnets.} \\
\\ 
\texttt{Let's analyze this based on my secret value function:} \\
\texttt{{-} **My values**: Remote: 0, Magnet: 1, Watch: 2.} \\
\texttt{{-} **Your offer**: You want the 2 remotes and 2 magnets.} \\
\texttt{{-} **My counter{-}offer logic**:} \\
\texttt{  {-} If I accept your proposal (you get remotes/magnets, I get watches):} \\
\texttt{    {-} You get: 2 Remotes (0 value) + 2 Magnets (1 value each) = **2 points**.} \\
\texttt{    {-} I get: 4 Watches (2 value each) = **8 points**.} \\
\texttt{  {-} This is a valid complementary split (2 remotes + 2 magnets + 4 watches = all items).} \\
\texttt{  {-} Since my goal is to maximize my score, getting 8 points is excellent.} \\
\\ 
\texttt{However, I need to check if there is a better outcome for *me* specifically, or if this deal is optimal.} \\
\texttt{{-} Total value of all items to me: (2 * 0) + (2 * 1) + (4 * 2) = 0 + 2 + 8 = 10.} \\
\texttt{{-} In this deal, I get 8 points. The other player gets 2 points.} \\
\texttt{{-} Is there a way for me to get more than 8? No, because the maximum possible value I can get from the watches is 8 (since there are only 4 watches). The remotes and magnets are worth 0 and 1 to me respectively. Even if I took the magnets, I would only gain 2 more points, but then I would have to give up some watches to the other player, which would reduce my score significantly (since watches are worth 2 each).} \\
\texttt{{-} Actually, let's re{-}evaluate: If I take the magnets (worth 1 each) and give up some watches, say I take 2 watches and 2 magnets. My score: (2*2) + (2*1) = 6. That's worse than 8.} \\
\texttt{{-} What if I take all items? I can't, because the other player must get something too, and if they get nothing, the deal fails (complementary requirement). But wait, the rule says "If your proposals are complementary...". It doesn't say both must get something, but typically in these games, if one player gets nothing, the other gets everything. However, the other player has their own values.} \\
\texttt{{-} Let's assume the other player's values are such that they value remotes/magnets highly. If I try to take everything, the other player might propose something that conflicts, leading to 0 for both.} \\
\texttt{{-} The safest and highest{-}scoring move for me is to} \\
            }
        }
    }
    & \\ \\

    \theutterance \stepcounter{utterance}  
    & \multicolumn{4}{p{0.6\linewidth}}{
        \cellcolor[rgb]{0.9,0.9,0.9}{
            \makecell[{{p{\linewidth}}}]{
                \texttt{\tiny{[P1$\rangle$GM]}}
                \texttt{I accept your offer. Since you value the watches highly and I value them at 0, this is a perfect trade for me. I will submit my secret proposal to take the 2 remotes and 2 magnets, while you take the 4 watches.} \\
\\ 
\texttt{{[}Proposal: 2 remote, 2 magnet{]}} \\
            }
        }
    }
    & & \\ \\

    \theutterance \stepcounter{utterance}  
    & & & \multicolumn{2}{p{0.3\linewidth}}{
        \cellcolor[rgb]{0.9,0.9,0.9}{
            \makecell[{{p{\linewidth}}}]{
                \texttt{\tiny{[GM$|$GM]}}
                \texttt{proposal} \\
            }
        }
    }
    & & \\ \\

    \theutterance \stepcounter{utterance}  
    & & & \multicolumn{2}{p{0.3\linewidth}}{
        \cellcolor[rgb]{0.9,0.9,0.9}{
            \makecell[{{p{\linewidth}}}]{
                \texttt{\tiny{[GM$|$GM]}}
                [2, 2, 0]
            }
        }
    }
    & & \\ \\

    \theutterance \stepcounter{utterance}  
    & & \multicolumn{4}{p{0.6\linewidth}}{
        \cellcolor[rgb]{0.9,0.9,0.9}{
            \makecell[{{p{\linewidth}}}]{
                \texttt{\tiny{[GM$\rangle$P2]}}
                \texttt{The other player has submitted a secret proposal. It is now your turn to submit a proposal of your own.} \\
            }
        }
    }
    & \\ \\

    \theutterance \stepcounter{utterance}  
    & & & \multicolumn{4}{p{0.6\linewidth}}{
        \cellcolor[rgb]{0.9,0.9,0.9}{
            \makecell[{{p{\linewidth}}}]{
                \texttt{\tiny{[GM$\langle$P2]}}
                \texttt{{[}Proposal: 4 watch{]}} \\
            }
        }
    }
     \\ \\

    \theutterance \stepcounter{utterance}  
    & & & \multicolumn{2}{p{0.3\linewidth}}{
        \cellcolor[rgb]{0.9,0.9,0.9}{
            \makecell[{{p{\linewidth}}}]{
                \texttt{\tiny{[GM$|$GM]}}
                \texttt{proposal} \\
            }
        }
    }
    & & \\ \\

    \theutterance \stepcounter{utterance}  
    & & & \multicolumn{2}{p{0.3\linewidth}}{
        \cellcolor[rgb]{0.9,0.9,0.9}{
            \makecell[{{p{\linewidth}}}]{
                \texttt{\tiny{[GM$|$GM]}}
                [0, 0, 4]
            }
        }
    }
    & & \\ \\

    \theutterance \stepcounter{utterance}  
    & & & \multicolumn{2}{p{0.3\linewidth}}{
        \cellcolor[rgb]{0.9,0.9,0.9}{
            \makecell[{{p{\linewidth}}}]{
                \texttt{\tiny{[GM$|$GM]}}
                [[2, 2, 0], [0, 0, 4]]
            }
        }
    }
    & & \\ \\

\end{supertabular}
}

\end{document}
